\section{Trabajos relacionados}

Convertir un comando en una estructura de intención y argumentos (slots) es
una tarea clásica de Natural Language Understanding (NLU). Weld et al.
\cite{weld2023survey} relevan los modelos conjuntos de detección de intención
y slot filling previos a los grandes modelos de lenguaje (LLMs), entrenados
específicamente para el dominio; son el punto de comparación natural del
enfoque actual, que resuelve la misma tarea generando directamente un JSON
con un modelo general. Bouraoui et al. \cite{bouraoui2018french} mostraron,
sobre un corpus de comandos de voz para hogar inteligente en francés, que los
recursos de NLU en inglés no transfieren bien a otras lenguas y que cada
idioma requiere corpus propios; este trabajo sigue esa lógica para el español
rioplatense, con un dataset de evaluación y no de entrenamiento.

El antecedente más cercano evalúa LLMs on-device para home assistants
\cite{ondevice2025homeassistant}: variantes cuantizadas de 16, 8 y 4 bits en
el doble rol de detección de intención y generación de respuesta, con
80\%--86\% de exactitud y latencias de 5 a 6 segundos. Este trabajo se
diferencia en el extremo inferior de tamaño (sub-2B, sin cuantización
agresiva), el idioma, y la interpretación estructurada en lugar de la
generación conversacional. El survey de Lu et al. \cite{gupta2025slms}
documenta que modelos de pocos miles de millones de parámetros alcanzan
paridad de tarea específica frente a modelos mucho mayores cuando se entrenan
con datos de alta calidad; los reportes técnicos de Qwen2.5
\cite{qwen2024report}, SmolLM2 \cite{allal2024smollm2} y OLMo 2
\cite{olmo2report} describen tres de las familias evaluadas aquí.

La decodificación restringida por gramática garantiza validez sintáctica con
independencia de la capacidad del modelo
\cite{beurerkellner2024guiding,dong2024xgrammar}. Este estudio optó
deliberadamente por no usarla, para medir el seguimiento de instrucciones en
un despliegue simple: resuelve la validez de formato, pero no evita que el
modelo alucine un valor semánticamente incorrecto dentro de un JSON válido,
patrón que la taxonomía de la Sección \ref{sec:dos-etapas} distingue.

La técnica central que este trabajo incorpora es la evaluación con un modelo
de lenguaje como juez. Zheng et al. \cite{zheng2023judging} muestran que un
LLM suficientemente capaz aproxima el juicio humano con acuerdo alto, y
proponen la técnica donde la coincidencia exacta subestima el desempeño real
— el caso de este trabajo, donde un sinónimo válido de un dispositivo no
coincide con el ground truth. Su limitación conocida es el sesgo de
auto-favorecimiento; como aquí el juez es uno de los modelos evaluados,
aplica directamente y se retoma en la Sección 6.
